\section{Auswertung}
\label{sec:Auswertung}
\input{build/results.tex}
\hypersetup{hidelinks}
\AtNextBibliography{\emergencystretch=2em}

Im Folgenden werden die Messergebnisse des Versuchs dokumentiert
und die Auswertung dieser beschrieben. \\

Die Messunsicherheit wird durch die Ablesegenauigkeit der verwendeten Messgeräte bestimmt.
Die Ablesegenauigkeit der Längenmessung beträgt $\pm\qty{1}{\milli\meter}$ und die der Winkelmessung 
$\pm\qty{1}{\degree}$. Die Ablesegenauigkeit der Leistungsmessung ist durch die Schwankung der gemessenen
Laserleistung begrenzt. Daher wird die Leistung bis zur letzten signifikanten Ziffer angegeben, die sich 
aus der Schwankung der Messwerte ergibt. Die Unsicherheit dieser Messwerte wird auf $\pm\num{1}$ der 
letzten signifikanten Ziffer geschätzt. Die Unsicherheiten der Messwerte werden in den Tabellen angegeben.
Die Fehlerfortpflanzung erfolgt mit dem Python-Paket \texttt{uncertainties} \cite{uncertainties}.

\subsection{Stabilitätsbedingung}
\label{subsec:v61-stabilitaet}

Für diesen Versuchsteil werden zwei Resonatoraufbauten mit unterschiedlichen Spiegelkonfigurationen untersucht.
Der konkav-konkave Resonator besteht aus zwei konkaven Spiegeln mit Krümmungsradien 
$r_1=r_2=\qty{1400}{\milli\meter}$, der plan-konkave Resonator aus einem konkaven Spiegel mit
Krümmungsradius $r_1=\qty{1400}{\milli\meter}$ und einem planen Spiegel. Im Folgenden wird die Notation $kk$
für die konkav-konkave und $pk$ für die plan-konkave Konfiguration verwendet.
Die Leistungen für die beiden Spiegelkonfigurationen sind in
\autoref{tab:v61-stabilitaet} zusammengestellt. Mit den angegebenen
Krümmungsradien ergeben sich die Stabilitätsparameter
\begin{align}
  G_{\mathrm{kk}}(L) &= g_1g_2
    =\left(1-\frac{L}{\qty{140}{\centi\meter}}\right)^2, \\
  G_{\mathrm{pk}}(L) &= g_1g_2
    =1-\frac{L}{\qty{140}{\centi\meter}}.
  \label{eq:v61-stabilitaet}
\end{align}
Dabei bezeichnet $L$ die Resonatorlänge und $g_1$ und $g_2$ die Stabilitätsparameter der beiden Spiegel.
Die Stabilitätsbedingung $0\leq g_1g_2\leq1$ liefert die oberen
geometrischen Grenzen $L_{\mathrm{kk,max}}=\qty{280}{\centi\meter}$ und
$L_{\mathrm{pk,max}}=\qty{140}{\centi\meter}$.
\autoref{fig:v61-stabilitaet} zeigt die Leistungen sowie die theoretischen
Stabilitätsparameter. Ein Leistungsfit wird nicht vorgenommen, da aus der
geometrischen Stabilitätsbedingung keine Leistungsfunktion folgt.

\input{build/table_stability.tex} 

\begin{figure}[H]
  \centering
  \includegraphics[width=0.85\textwidth]{build/stabilitaet.pdf}
  \caption{Gemessene Leistung und theoretische Stabilitätsparameter der beiden Resonatoren.
  Der grüne Bereich kennzeichnet $0\leq g_1g_2\leq1$, die senkrechten
  Linien markieren die theoretischen oberen Längengrenzen.}
  \label{fig:v61-stabilitaet}
\end{figure}
\FloatBarrier

\subsection{Transversale Moden}
\label{subsec:v61-tem}

Die Leistungsprofile werden mit Fits an die Hermite-Gauß-Verteilungen
unter Berücksichtigung eines konstanten Untergrunds $B$ ausgewertet:
\begin{align}
  P_{00}(x) &= B+A\exp\left[-2\frac{(x-x_0)^2}{w^2}\right],
    \label{eq:v61-tem00} \\
  P_{10}(x) &= B+A\,2\mathrm{e}\,\frac{(x-x_0)^2}{w^2}
    \exp\left[-2\frac{(x-x_0)^2}{w^2}\right].
    \label{eq:v61-tem10}
\end{align}
Dabei ist $w$ ein Maß für die Breite des Strahls und $x_0$ die Profilmitte.
Die Normierung ist so gewählt, dass $A$ bei beiden Modellen die
Maximalleistung über dem Untergrund angibt.

\subsubsection{Grundmode TEM\texorpdfstring{$_{00}$}{00}}

\autoref{tab:v61-tem00} enthält sämtliche Messwerte der Grundmode.
Der Fit in \autoref{fig:v61-tem00} ergibt
\begin{align*}
  A_{00}&=\vZeroA, & x_{0,00}&=\vZeroCentre, \\
  w_{00}&=\vZeroWidth, & B_{00}&=\vZeroB.
\end{align*}
Mit $\nu=\vZeroDof$ Freiheitsgraden beträgt das reduzierte
$\chi^2_\nu=\vZeroChi$.

\input{build/table_tem00.tex}
\begin{figure}[H]
  \centering
  \begin{subfigure}{0.85\textwidth}
    \centering
    \includegraphics[width=\linewidth]{build/tem00.pdf}
    \caption{Gemessenes Leistungsprofil mit Fit.}
  \end{subfigure}

  \medskip
  \begin{subfigure}{0.75\textwidth}
    \centering
    \includegraphics[trim={1410bp 1169bp 1722bp 1180bp},clip,width=\linewidth]{build/tem00_foto.png}
    \caption{Ausschnitt der fotografischen Aufnahme der TEM$_{00}$-Mode.}
  \end{subfigure}
  \caption{Gemessenes TEM$_{00}$-Profil mit Fit nach
  \autoref{eq:v61-tem00} und fotografische Aufnahme des Modenbildes.}
  \label{fig:v61-tem00}
\end{figure}
\FloatBarrier

\subsubsection{Erste transversale Mode TEM\texorpdfstring{$_{10}$}{10}}

Die Messwerte in \autoref{tab:v61-tem10} zeigen zwei Maxima mit einem
Minimum dazwischen, charakteristisch für die TEM$_{10}$-Mode. Der Fit der Messwerte
an \autoref{eq:v61-tem10} ergeben
\begin{align*}
  A_{10}&=\vFirstA, & x_{0,10}&=\vFirstCentre, \\
  w_{10}&=\vFirstWidth, & B_{10}&=\vFirstB.
\end{align*}
Das reduzierte $\chi^2_\nu=\vFirstChi$ gehört zu
$\nu=\vFirstDof$ Freiheitsgraden.
\autoref{fig:v61-tem10} zeigt den zugehörigen Fit und die fotografische Aufnahme.

\input{build/table_tem10.tex}
\begin{figure}[H]
  \centering
  \begin{subfigure}{0.85\textwidth}
    \centering
    \includegraphics[width=\linewidth]{build/tem10.pdf}
    \caption{Gemessenes Leistungsprofil mit Fit.}
  \end{subfigure}

  \medskip
  \begin{subfigure}{0.75\textwidth}
    \centering
    \includegraphics[trim={1330bp 1219bp 1582bp 965bp},clip,width=\linewidth]{build/tem10_foto.png}
    \caption{Ausschnitt der fotografischen Aufnahme der TEM$_{10}$-Mode.}
  \end{subfigure}
  \caption{Gemessenes TEM$_{10}$-Profil mit Fit an die symmetrische
  Hermite-Gauß-Verteilung nach \autoref{eq:v61-tem10} und fotografische
  Aufnahme des Modenbildes.}
  \label{fig:v61-tem10}
\end{figure}
\FloatBarrier

\subsection{Polarisation}
\label{subsec:v61-polarisation}

Die Winkelabhängigkeit der Leistung aus
\autoref{tab:v61-polarisation} wird mit einem Fit an das Malus-Gesetz
einschließlich eines konstanten Untergrundterms ausgewertet:
\begin{equation}
  P(\varphi)=B+A\cos^2(\varphi-\varphi_0).
  \label{eq:v61-malus}
\end{equation}
Der Fit in \autoref{fig:v61-polarisation} liefert
\begin{equation*}
  A=\vPolA,\qquad \varphi_0=\vPolTheta,\qquad B=\vPolB.
\end{equation*}
Die Fitgüte beträgt $\chi^2_\nu=\vPolChi$ bei $\nu=\vPolDof$.

\input{build/table_polarisation.tex}
\begin{figure}[H]
  \centering
  \includegraphics[width=0.85\textwidth]{build/polarisation.pdf}
  \caption{Leistung als Funktion der Analysatorstellung und Fit
  des Malus-Gesetzes. Winkel- und Leistungsunsicherheiten gehen in die
  Fitberechnung ein.}
  \label{fig:v61-polarisation}
\end{figure}
\FloatBarrier

\subsection{Multimodenbetrieb und Frequenzspektrum}
\label{subsec:v61-longitudinal}

Die Peaks werden im Frequenzspektrum abgelesen, das mit Hilfe einer
schnellen Fouriertransformation (FFT) der longitudinalen Modenschwebung
gewonnen wird. Ihre Frequenzpositionen und Abstände werden mit dem
theoretisch vorhergesagten Modenabstand $\Delta\nu$ verglichen. Es gilt
\begin{equation}
  \Delta\nu=\frac{c}{2L},\qquad f_n(L)=n\frac{c}{2L},\qquad
  n=1,2,\ldots.
  \label{eq:v61-longitudinal}
\end{equation}
Die gemessenen Peakpositionen sind in \autoref{tab:v61-longitudinal}
angegeben. \autoref{tab:v61-mode-spacings} stellt die daraus ermittelten
Modenabstände $\overline{\Delta\nu}$ den theoretischen Werten
$\Delta\nu_{\mathrm{theo}}=c/(2L)$ gegenüber.

\input{build/table_longitudinal.tex}
\input{build/table_mode_spacings.tex}
\begin{figure}[H]
  \centering
  \includegraphics[width=0.85\textwidth]{build/longitudinalmoden.pdf}
  \caption{Aus den Peakpositionen ermittelte Modenabstände als Funktion
  der Resonatorlänge und theoretische Erwartung $c/(2L)$.
  Unten ist die Differenz $d_\nu$ dargestellt.}
  \label{fig:v61-longitudinal}
\end{figure}

Für die Abschätzung der Dopplerbreite der Neon-Emission wird
mangels gemessener Gastemperatur
$T=\qty{293.15}{\kelvin}$ angenommen.

Mit \mbox{$m_{\mathrm{Ne}}=\qty{20.1797}{u}$ \cite{NIST_Neon}}
und \mbox{$\lambda_{\mathrm{ref}}=\qty{632.8}{\nano\meter}$ \cite[2]{eLas_HeNe}} ergibt sich die
Halbwertsbreite
\begin{equation}
  \Delta\nu_{\mathrm D}
  =\frac{1}{\lambda_{\mathrm{ref}}}
    \sqrt{\frac{8 k_{\mathrm B}T\ln 2}{m_{\mathrm{Ne}}}}
  \approx\vDoppler\,\unit{\giga\hertz}.
  \label{eq:v61-doppler}
\end{equation}

\subsection{Wellenlänge}
\label{subsec:v61-wellenlaenge}

Die Gitter besitzen die Liniendichten $N_1=\qty{80}{\per\milli\meter}$
und $N_2=\qty{100}{\per\milli\meter}$. Mit der Gitterkonstante $b=1/N$
gilt bei senkrechtem Einfall
\begin{equation}
  m\lambda=b\sin\theta_m.
  \label{eq:v61-gitterbedingung}
\end{equation}
Für die Position $y_m$ des Maximums der Ordnung $m$ folgt aus der
Schirmgeometrie mit dem Schirmabstand $d$ und der Position $y_0$
des zentralen Maximums
\begin{equation}
  y_m-y_0=d\tan\theta_m
  =d\frac{\sin\theta_m}{\sqrt{1-\sin^2\theta_m}}.
\end{equation}
Einsetzen von $\sin\theta_m=m\lambda/b$ und Umstellen nach $y_m$
liefert die Fitfunktion
\begin{equation}
  y_m=y_0+d\frac{m\lambda/b}{\sqrt{1-(m\lambda/b)^2}}.
  \label{eq:v61-gitter}
\end{equation}
Die Messwerte in \autoref{tab:v61-wellenlaenge} werden in der notierten
Reihenfolge den Ordnungen $m=2,1,0,-1,-2$ zugeordnet und für jedes Gitter
getrennt gefittet. Dabei werden die Positions- und Schirmabstandsunsicherheiten
berücksichtigt und $y_0$ jeweils mitbestimmt.

\input{build/table_wavelength.tex}

Die Fits in \autoref{fig:v61-wellenlaenge} liefern
\begin{align}
  N_1=\qty{80}{\per\milli\meter}:\quad
  &\lambda_1=\vWaveOne,\quad
  \chi^2_{\nu,1}=\vWaveOneChi,\quad \nu_1=\vWaveOneDof, \\
  N_2=\qty{100}{\per\milli\meter}:\quad
  &\lambda_2=\vWaveTwo,\quad
  \chi^2_{\nu,2}=\vWaveTwoChi,\quad \nu_2=\vWaveTwoDof.
  \label{eq:v61-wellenlaenge-ergebnis}
\end{align}

\begin{figure}[H]
  \centering
  \includegraphics[width=\textwidth]{build/wellenlaenge.pdf}
  \caption{Messwerte und getrennte Fits nach \autoref{eq:v61-gitter}
  für die beiden Gitter. Die gestrichelten Kurven zeigen die Erwartung
  für \qty{632.8}{\nano\meter} bei den jeweils fitbestimmten
  Positionsnullpunkten und notierten Schirmabständen.}
  \label{fig:v61-wellenlaenge}
\end{figure}
\FloatBarrier
